\documentclass{amsart}
% For dealing with references we use the comment environment
\usepackage{verbatim}
\newenvironment{reference}{\comment}{\endcomment}
%\newenvironment{reference}{}{}
\newenvironment{slogan}{\comment}{\endcomment}
\newenvironment{history}{\comment}{\endcomment}

% For commutative diagrams we use Xy-pic
\usepackage[all]{xy}

% We use 2cell for 2-commutative diagrams.
\xyoption{2cell}
\UseAllTwocells

% We use multicol for the list of chapters between chapters
\usepackage{multicol}

% This is generally recommended for better output
\usepackage{lmodern}
\usepackage[T1]{fontenc}

% For cross-file-references

% Package for hypertext links:
\usepackage{hyperref}

% For any local file, say "hello.tex" you want to link to please

% Theorem environments.
%
\theoremstyle{plain}
\newtheorem{theorem}[subsection]{Theorem}
\newtheorem{proposition}[subsection]{Proposition}
\newtheorem{lemma}[subsection]{Lemma}

\theoremstyle{definition}
\newtheorem{definition}[subsection]{Definition}
\newtheorem{example}[subsection]{Example}
\newtheorem{exercise}[subsection]{Exercise}
\newtheorem{situation}[subsection]{Situation}

\theoremstyle{remark}
\newtheorem{remark}[subsection]{Remark}
\newtheorem{remarks}[subsection]{Remarks}

\numberwithin{equation}{subsection}

% Macros
%
\def\lim{\mathop{\mathrm{lim}}\nolimits}
\def\colim{\mathop{\mathrm{colim}}\nolimits}
\def\Spec{\mathop{\mathrm{Spec}}}
\def\Hom{\mathop{\mathrm{Hom}}\nolimits}
\def\Ext{\mathop{\mathrm{Ext}}\nolimits}
\def\SheafHom{\mathop{\mathcal{H}\!\mathit{om}}\nolimits}
\def\SheafExt{\mathop{\mathcal{E}\!\mathit{xt}}\nolimits}
\def\Sch{\mathit{Sch}}
\def\Mor{\mathop{\mathrm{Mor}}\nolimits}
\def\Ob{\mathop{\mathrm{Ob}}\nolimits}
\def\Sh{\mathop{\mathit{Sh}}\nolimits}
\def\NL{\mathop{N\!L}\nolimits}
\def\CH{\mathop{\mathrm{CH}}\nolimits}
\def\proetale{{pro\text{-}\acute{e}tale}}
\def\etale{{\acute{e}tale}}
\def\QCoh{\mathit{QCoh}}
\def\Ker{\mathop{\mathrm{Ker}}}
\def\Im{\mathop{\mathrm{Im}}}
\def\Coker{\mathop{\mathrm{Coker}}}
\def\Coim{\mathop{\mathrm{Coim}}}

% Boxtimes
%
\DeclareMathSymbol{\boxtimes}{\mathbin}{AMSa}{"02}

%
% Macros for moduli stacks/spaces
%
\def\QCohstack{\mathcal{QC}\!\mathit{oh}}
\def\Cohstack{\mathcal{C}\!\mathit{oh}}
\def\Spacesstack{\mathcal{S}\!\mathit{paces}}
\def\Quotfunctor{\mathrm{Quot}}
\def\Hilbfunctor{\mathrm{Hilb}}
\def\Curvesstack{\mathcal{C}\!\mathit{urves}}
\def\Polarizedstack{\mathcal{P}\!\mathit{olarized}}
\def\Complexesstack{\mathcal{C}\!\mathit{omplexes}}
% \Pic is the operator that assigns to X its picard group, usage \Pic(X)
% \Picardstack_{X/B} denotes the Picard stack of X over B
% \Picardfunctor_{X/B} denotes the Picard functor of X over B
\def\Pic{\mathop{\mathrm{Pic}}\nolimits}
\def\Picardstack{\mathcal{P}\!\mathit{ic}}
\def\Picardfunctor{\mathrm{Pic}}
\def\Deformationcategory{\mathcal{D}\!\mathit{ef}}

\setlength{\emergencystretch}{3em}
\begin{document}
\title[Stacks Project, Tag 0BHU]{737 --- Regular alterations of local Noetherian domains force normal completed normalization and regular generic formal fibre}
\maketitle
\noindent Copyright (C) 2005--2025 Johan de Jong. Stacks Project authors. Source: \href{https://stacks.math.columbia.edu/tag/0BHU}{Tag 0BHU}.

\noindent Editorial difficulty note (not author text). Difficulty H2: The author passes through the finite algebra of global sections of the alteration, separates its semilocal completions, proves normality through proper regular base change, and descends regularity of the generic formal fibre. This interaction of alterations, normalization and potentially bad formal fibres is above qualification commutative algebra..

\noindent Editorial provenance note (not author text). History: excerpt from revision \texttt{a0444\allowbreak 6e57e\allowbreak c1fbc\allowbreak 252a8\allowbreak 71afc\allowbreak ec775\allowbreak 2fb28\allowbreak 07b14}, spaces-resolve.tex, lines 753--831. Reference presentation was adapted; original-author context and core support blocks were retained or added. The mathematical wording is unchanged. Licensed under GNU FDL 1.2 or later; no invariant sections or cover texts. See ../licenses/COPYING. Context blocks reproduce author definitions and supporting results; they are not additional counted problems.

\begin{definition}
\label{spaces-over-fields-definition-alteration}
Let $S$ be a scheme. Let $X$ be an integral algebraic space over $S$.
An {\it alteration of $X$} is a proper dominant morphism $f : Y \to X$
of algebraic spaces over $S$ with $Y$ integral such that $f^{-1}(U) \to U$
is finite for some nonempty open $U \subset X$.
\end{definition}

\begin{lemma}
\label{lemma-regular-alteration-implies}
Let $S$ be a scheme. Let $Y$ be a Noetherian integral algebraic space
over $S$. Assume there exists an alteration
$f : X \to Y$ with $X$ regular. Then the normalization $Y^\nu \to Y$
is finite and $Y$ has a dense open which is regular.
\end{lemma}

\begin{proof}
By \'etale localization, it suffices to prove this when
$Y = \Spec(A)$ where $A$ is a Noetherian domain.
Let $B$ be the integral closure of $A$ in its fraction field.
Set $C = \Gamma(X, \mathcal{O}_X)$. By
Cohomology of Spaces, Lemma
\href{https://stacks.math.columbia.edu/tag/08AR}{lemma proper pushforward coherent (Tag 08AR)}
we see that $C$ is a finite $A$-module.
As $X$ is normal
(Properties of Spaces, Lemma
\href{https://stacks.math.columbia.edu/tag/0BGT}{lemma regular normal (Tag 0BGT)})
we see that $C$ is normal domain
(Spaces over Fields, Lemma
\href{https://stacks.math.columbia.edu/tag/0BH3}{lemma normal integral sections (Tag 0BH3)}).
Thus $B \subset C$ and we conclude that $B$ is finite over $A$
as $A$ is Noetherian.

\medskip\noindent
There exists a nonempty open $V \subset Y$ such that $f^{-1}V \to V$
is finite, see Spaces over Fields, Definition
\ref{spaces-over-fields-definition-alteration}.
After shrinking $V$ we may assume that $f^{-1}V \to V$ is flat
(Morphisms of Spaces, Proposition
\href{https://stacks.math.columbia.edu/tag/06QS}{proposition generic flatness reduced (Tag 06QS)}).
Thus $f^{-1}V \to V$ is faithfully flat. Then $V$ is regular by
Algebra, Lemma \href{https://stacks.math.columbia.edu/tag/07NG}{lemma descent regular (Tag 07NG)}.
\end{proof}


\begin{lemma}
\label{lemma-port-regularity-to-completion}
Let $(A, \mathfrak m, \kappa)$ be a Noetherian local ring.
Let $X \to \Spec(A)$ be a morphism which is locally of finite type
with $X$ a decent algebraic space. Set
$Y = X \times_{\Spec(A)} \Spec(A^\wedge)$. Let $y \in |Y|$
with image $x \in |X|$. Then
\begin{enumerate}
\item if $\mathcal{O}_{Y, y}^h$ is regular, then
$\mathcal{O}_{X, x}^h$ is regular,
\item if $y$ is in the closed fibre, then $\mathcal{O}_{Y, y}^h$ is regular
$\Leftrightarrow \mathcal{O}_{X, x}^h$ is regular, and
\item If $X$ is proper over $A$, then $X$ is regular
if and only if $Y$ is regular.
\end{enumerate}
\end{lemma}

\begin{proof}
By \'etale localization the first two statements follow
immediately from the counter part to this lemma for schemes, see
Resolution of Surfaces, Lemma \ref{resolve-lemma-port-regularity-to-completion}.
For part (3), since $Y \to X$ is surjective (as $A \to A^\wedge$
is faithfully flat) we see that $Y$ regular implies $X$ regular
by part (1). Conversely, if $X$ is regular, then the henselian
local rings of $Y$ are regular for all points of the special fibre.
Let $y \in |Y|$ be a general point.
Since $|Y| \to |\Spec(A^\wedge)|$ is closed in the proper
case, we can find a specialization $y \leadsto y_0$ with
$y_0$ in the closed fibre. Choose an elementary \'etale
neighbourhood $(V, v_0) \to (Y, y_0)$ as in
Decent Spaces, Lemma
\href{https://stacks.math.columbia.edu/tag/0BBP}{lemma decent space elementary etale neighbourhood (Tag 0BBP)}.
Since $Y$ is decent we can lift $y \leadsto y_0$ to a specialization
$v \leadsto v_0$ in $V$
(Decent Spaces, Lemma \href{https://stacks.math.columbia.edu/tag/03IL}{lemma decent specialization (Tag 03IL)}).
Then we conclude that
$\mathcal{O}_{V, v}$ is a localization of $\mathcal{O}_{V, v_0}$
hence regular and the proof is complete.
\end{proof}


\begin{lemma}
\label{resolve-lemma-algebra-helper}
Let $(A, \mathfrak m)$ be a local Noetherian ring. Let $B \subset C$
be finite $A$-algebras. Assume that (a) $B$ is a normal ring, and
(b) the $\mathfrak m$-adic completion $C^\wedge$ is a normal ring.
Then $B^\wedge$ is a normal ring.
\end{lemma}

\begin{proof}
Consider the commutative diagram
$$
\xymatrix{
B \ar[r] \ar[d] & C \ar[d] \\
B^\wedge \ar[r] & C^\wedge
}
$$
Recall that $\mathfrak m$-adic completion on the category of
finite $A$-modules is exact because it is given by tensoring with
the flat $A$-algebra $A^\wedge$
(Algebra, Lemma \href{https://stacks.math.columbia.edu/tag/00MB}{lemma completion flat (Tag 00MB)}).
We will use Serre's criterion
(Algebra, Lemma \href{https://stacks.math.columbia.edu/tag/031S}{lemma criterion normal (Tag 031S)})
to prove that the Noetherian ring $B^\wedge$ is normal.
Let $\mathfrak q \subset B^\wedge$ be a prime lying over
$\mathfrak p \subset B$. If $\dim(B_\mathfrak p) \geq 2$, then
$\text{depth}(B_\mathfrak p) \geq 2$ and since
$B_\mathfrak p \to B^\wedge_\mathfrak q$ is flat we find
that $\text{depth}(B^\wedge_\mathfrak q) \geq 2$
(Algebra, Lemma \href{https://stacks.math.columbia.edu/tag/0337}{lemma apply grothendieck (Tag 0337)}).
If $\dim(B_\mathfrak p) \leq 1$, then $B_\mathfrak p$ is
either a discrete valuation ring or a field.
In that case $C_\mathfrak p$ is faithfully flat over $B_\mathfrak p$
(because it is finite and torsion free).
Hence $B^\wedge_\mathfrak p \to C^\wedge_\mathfrak p$ is
faithfully flat and the same holds after localizing at $\mathfrak q$.
As $C^\wedge$ and hence any localization is $(S_2)$ we conclude that
$B^\wedge_\mathfrak p$ is $(S_2)$ by
Algebra, Lemma \href{https://stacks.math.columbia.edu/tag/0352}{lemma descent Sk (Tag 0352)}.
All in all we find that
$(S_2)$ holds for $B^\wedge$. To prove that $B^\wedge$ is
$(R_1)$ we only have to consider primes $\mathfrak q \subset B^\wedge$
with $\dim(B^\wedge_\mathfrak q) \leq 1$. Since
$\dim(B^\wedge_\mathfrak q) = \dim(B_\mathfrak p) +
\dim(B^\wedge_\mathfrak q/\mathfrak p B^\wedge_\mathfrak q)$ by
Algebra, Lemma \href{https://stacks.math.columbia.edu/tag/00OM}{lemma dimension base fibre total (Tag 00OM)}
we find that $\dim(B_\mathfrak p) \leq 1$ and
we see that $B^\wedge_\mathfrak q \to C^\wedge_\mathfrak q$
is faithfully flat as before. We conclude using
Algebra, Lemma \href{https://stacks.math.columbia.edu/tag/0353}{lemma descent Rk (Tag 0353)}.
\end{proof}


\begin{lemma}
\label{resolve-lemma-port-regularity-to-completion}
Let $(A, \mathfrak m, \kappa)$ be a Noetherian local ring.
Let $X \to \Spec(A)$ be a morphism which is locally of finite type.
Set $Y = X \times_{\Spec(A)} \Spec(A^\wedge)$. Let $y \in Y$ with
image $x \in X$. Then
\begin{enumerate}
\item if $\mathcal{O}_{Y, y}$ is regular, then $\mathcal{O}_{X, x}$
is regular,
\item if $y$ is in the closed fibre, then $\mathcal{O}_{Y, y}$ is regular
$\Leftrightarrow \mathcal{O}_{X, x}$ is regular, and
\item If $X$ is proper over $A$, then $X$ is regular
if and only if $Y$ is regular.
\end{enumerate}
\end{lemma}

\begin{proof}
Since $A \to A^\wedge$ is faithfully flat
(Algebra, Lemma \href{https://stacks.math.columbia.edu/tag/00MC}{lemma completion faithfully flat (Tag 00MC)}),
we see that $Y \to X$ is flat. Hence (1) by
Algebra, Lemma \href{https://stacks.math.columbia.edu/tag/07NG}{lemma descent regular (Tag 07NG)}.
Lemma \ref{resolve-lemma-iso-completions} shows the morphism $Y \to X$
induces an isomorphism on complete local rings at points
of the special fibres. Thus (2) by
More on Algebra, Lemma \href{https://stacks.math.columbia.edu/tag/07NY}{lemma completion regular (Tag 07NY)}.
If $X$ is proper over $A$, then $Y$ is proper over $A^\wedge$
(Morphisms, Lemma \href{https://stacks.math.columbia.edu/tag/01W4}{lemma base change proper (Tag 01W4)})
and we see every closed point of $X$ and $Y$ lies in the closed fibre.
Thus we see that $Y$ is a regular scheme if and only if $X$ is so by
Properties, Lemma \href{https://stacks.math.columbia.edu/tag/02IT}{lemma characterize regular (Tag 02IT)}.
\end{proof}


\begin{lemma}
\label{resolve-lemma-iso-completions}
Let $(A, \mathfrak m, \kappa)$ be a local ring with finitely generated
maximal ideal $\mathfrak m$. Let $X$ be a scheme over $A$.
Let $Y = X \times_{\Spec(A)} \Spec(A^\wedge)$ where
$A^\wedge$ is the $\mathfrak m$-adic completion of $A$.
For a point $q \in Y$ with image $p \in X$ lying
over the closed point of $\Spec(A)$ the
local ring map $\mathcal{O}_{X, p} \to \mathcal{O}_{Y, q}$
induces an isomorphism on completions.
\end{lemma}

\begin{proof}
We may assume $X$ is affine. Then we may write $X = \Spec(B)$.
Let $\mathfrak q \subset B' = B \otimes_A A^\wedge$ be the
prime corresponding to $q$ and let $\mathfrak p \subset B$
be the prime ideal corresponding to $p$.
By Algebra, Lemma \href{https://stacks.math.columbia.edu/tag/05GG}{lemma hathat finitely generated (Tag 05GG)}
we have
$$
B'/(\mathfrak m^\wedge)^n B' =
A^\wedge/(\mathfrak m^\wedge)^n \otimes_A B =
A/\mathfrak m^n \otimes_A B = B/\mathfrak m^n B
$$
for all $n$. Since $\mathfrak m B \subset \mathfrak p$ and
$\mathfrak m^\wedge B' \subset \mathfrak q$ we see that
$B/\mathfrak p^n$ and $B'/\mathfrak q^n$ are both
quotients of the ring displayed above by the $n$th power
of the same prime ideal. The lemma follows.
\end{proof}


\begin{lemma}
\label{lemma-regular-alteration-implies-local}
Let $(A, \mathfrak m, \kappa)$ be a local Noetherian domain.
Assume there exists an alteration $f : X \to \Spec(A)$
with $X$ regular. Then
\begin{enumerate}
\item there exists a nonzero $f \in A$ such that $A_f$ is regular,
\item the integral closure $B$ of $A$ in its fraction field is finite over $A$,
\item the $\mathfrak m$-adic completion of $B$ is a normal ring, i.e., the
completions of $B$ at its maximal ideals are normal domains, and
\item the generic formal fibre of $A$ is regular.
\end{enumerate}
\end{lemma}

\begin{proof}
Parts (1) and (2) follow from Lemma \ref{lemma-regular-alteration-implies}.
We have to redo part of the proof of that lemma in order to set up notation
for the proof of (3). Set $C = \Gamma(X, \mathcal{O}_X)$. By
Cohomology of Spaces, Lemma
\href{https://stacks.math.columbia.edu/tag/08AR}{lemma proper pushforward coherent (Tag 08AR)}
we see that $C$ is a finite $A$-module.
As $X$ is normal
(Properties of Spaces, Lemma
\href{https://stacks.math.columbia.edu/tag/0BGT}{lemma regular normal (Tag 0BGT)})
we see that $C$ is normal domain
(Spaces over Fields, Lemma
\href{https://stacks.math.columbia.edu/tag/0BH3}{lemma normal integral sections (Tag 0BH3)}).
Thus $B \subset C$ and we conclude that $B$ is finite over $A$
as $A$ is Noetherian. By
Resolution of Surfaces, Lemma \ref{resolve-lemma-algebra-helper}
in order to prove (3) it suffices to show
that the $\mathfrak m$-adic completion $C^\wedge$ is normal.

\medskip\noindent
By Algebra, Lemma \href{https://stacks.math.columbia.edu/tag/07N9}{lemma completion finite extension (Tag 07N9)}
the completion $C^\wedge$ is the product of the completions of
$C$ at the prime ideals of $C$ lying over $\mathfrak m$.
There are finitely many of these and these are the maximal
ideals $\mathfrak m_1, \ldots, \mathfrak m_r$ of $C$.
(The corresponding result for $B$ explains the final statement of the lemma.)
Thus replacing $A$ by $C_{\mathfrak m_i}$ and $X$ by
$X_i = X \times_{\Spec(C)} \Spec(C_{\mathfrak m_i})$
we reduce to the case discussed in the next paragraph.
(Note that $\Gamma(X_i, \mathcal{O}) = C_{\mathfrak m_i}$ by
Cohomology of Spaces,
Lemma \href{https://stacks.math.columbia.edu/tag/073K}{lemma flat base change cohomology (Tag 073K)}.)

\medskip\noindent
Here $A$ is a Noetherian local normal domain and $f : X \to \Spec(A)$
is a regular alteration with $\Gamma(X, \mathcal{O}_X) = A$.
We have to show that the completion $A^\wedge$
of $A$ is a normal domain. By
Lemma \ref{lemma-port-regularity-to-completion}
$Y = X \times_{\Spec(A)} \Spec(A^\wedge)$ is regular.
Since $\Gamma(Y, \mathcal{O}_Y) = A^\wedge$
by Cohomology of Spaces,
Lemma \href{https://stacks.math.columbia.edu/tag/073K}{lemma flat base change cohomology (Tag 073K)}.
We conclude that $A^\wedge$ is normal as before.
Namely, $Y$ is normal by Properties of Spaces, Lemma
\href{https://stacks.math.columbia.edu/tag/0BGT}{lemma regular normal (Tag 0BGT)}.
It is connected because $\Gamma(Y, \mathcal{O}_Y) = A^\wedge$ is local.
Hence $Y$ is normal and integral (as connected and normal
implies integral for separated algebraic spaces). Thus
$\Gamma(Y, \mathcal{O}_Y) = A^\wedge$ is a normal domain by
Spaces over Fields, Lemma
\href{https://stacks.math.columbia.edu/tag/0BH3}{lemma normal integral sections (Tag 0BH3)}.
This proves (3).

\medskip\noindent
Proof of (4). Let $\eta \in \Spec(A)$ denote the generic point
and denote by a subscript $\eta$ the base change to $\eta$.
Since $f$ is an alteration, the scheme $X_\eta$ is finite and
faithfully flat over $\eta$. Since $Y = X \times_{\Spec(A)} \Spec(A^\wedge)$
is regular by Lemma \ref{lemma-port-regularity-to-completion}
we see that $Y_\eta$ is regular (as a limit of opens in $Y$).
Then $Y_\eta \to \Spec(A^\wedge \otimes_A \kappa(\eta))$ is finite
faithfully flat onto the generic formal fibre. We conclude by
Algebra, Lemma \href{https://stacks.math.columbia.edu/tag/07NG}{lemma descent regular (Tag 07NG)}.
\end{proof}
\end{document}
